\section*{第 \ref{ch:b1:counting} 章——计数}

\begin{solution}{exo:b1:counting:1}
各步相互独立，由乘法原理：$26^2 \times 10^3 \times 26^2
= 26^4 \times 1000 = 456\,976\,000$ 个车牌。若四个字母
两两不同，则字母部分构成字母表的 $4$-排列，
共 $26 \times 25 \times 24 \times 23 = 358\,800$ 种，
故 $358\,800 \times 1000 = 358\,800\,000$ 个车牌。
\end{solution}

\begin{solution}{exo:b1:counting:2}
\textsc{orange} 有 $6$ 个互不相同的字母：$6! = 720$ 种字母重排。
\textsc{banana} 有 $6$ 个字母，其中有重复（$3$ 个 a，$2$ 个 n，
$1$ 个 b）：每种重排由 a 所在的位置
（$\binom 63$ 种选法）确定，再在其余 $3$ 个位置中确定 n 的位置，
（$\binom 32$ 种选法），b 占最后一个位置：
$\binom{6}{3}\binom{3}{2} = 20 \times 3 = 60$ 种重排
（等价地，$6!/(3!\,2!\,1!) = 60$）。
\end{solution}

\begin{solution}{exo:b1:counting:3}
总数：$\binom{12}{4} = 495$。恰有 $2$ 名女性：先选女性
（$\binom 72 = 21$），再选 $2$ 名男性（$\binom 52 = 10$）：$210$ 个委员会。
至少一名男性：取其补集“没有男性”，
故 $\binom{12}{4} - \binom{7}{4} = 495 - 35 = 460$。
\end{solution}

\begin{solution}{exo:b1:counting:4}
\emph{生日：}以 $12$ 个月份为抽屉，$13$ 人放入 $12$
个抽屉，必有两人落在同一抽屉
（见\cref{cor:b1:counting:pigeonhole}）。

\emph{连续整数：}以 $n$ 个\omterm{def:b1:logic:sets}{集合} $\{1,2\},
\{3,4\}, \dots, \{2n-1, 2n\}$，构成 $\intint{1}{2n}$ 的一个\omterm{thm:b1:logic:partition}{划分}。
任取 $n + 1$ 个整数，必有两个落在同一\omterm{def:b1:logic:sets}{集合}中，
而同一\omterm{def:b1:logic:sets}{集合}中的两个元素必为相邻整数。
\end{solution}

\begin{solution}{exo:b1:counting:5}
对 $1 \leq k \leq n$，
\[
k \binom nk = k\,\frac{n!}{k!\,(n-k)!}
= n\,\frac{(n-1)!}{(k-1)!\,(n-k)!} = n \binom{n-1}{k-1}.
\]
求和并作指标代换 $j = k - 1$，得
\[
\sum_{k=0}^{n} k \binom nk = n \sum_{j=0}^{n-1} \binom{n-1}{j}
= n\, 2^{n-1}
\]
由\cref{prop:b1:counting:identities}可得。（另证：
对 $(1+x)^n = \sum_k \binom nk x^k$ 求导，并令 $x = 1$。）
\end{solution}

\begin{solution}{exo:b1:counting:6}
严格递增\omterm{def:b1:logic:map}{映射} $f \colon \intint{1}{k} \to \intint{1}{n}$ 由其像
完全确定，其像为一个 $k$ 元子集（即 $\intint{1}{n}$ 的子集，按递增顺序列出）；
反之，每个 $k$ 元子集恰给出一个这样的\omterm{def:b1:logic:map}{映射}。
于是严格递增\omterm{def:b1:logic:map}{映射}共 $\binom nk$ 个。

若 $f$ 只是递增的，令 $g(i) = f(i) + i - 1$。则 $g$ 是
严格递增的（相邻自变量之间，$f$ 的增量 $\geq 0$，
$i - 1$ 的增量为 $1$），取值于 $\intint{1}{n + k - 1}$；
$f(i) = g(i) - i + 1$ 还原 $f$，其中 $g$ 为任一严格递增的
映入 $\intint{1}{n+k-1}$ 的\omterm{def:b1:logic:map}{映射}。这是\omterm{def:b1:logic:inj}{双射}。
故共有 $\binom{n + k - 1}{k}$ 个递增\omterm{def:b1:logic:map}{映射}。
\end{solution}

\begin{solution}{exo:b1:counting:7}
将\omterm{def:b1:logic:sets}{集合} $E$（含 $m + n$ 个元素）分为两块 $M$（含 $m$ 个元素）
和 $N$（含 $n$ 个元素）。一个 $k$ 元子集（$E$ 的子集）含 $j$ 个元素
取自 $M$（$0 \leq j \leq k$），其余 $k - j$ 个取自 $N$；对固定的 $j$，这样的子集
共 $\binom mj \binom{n}{k-j}$ 个，而 $j = 0, \dots,
k$ 的各情形恰好\omterm{thm:b1:logic:partition}{划分}全部 $k$ 元子集。由加法原理得
范德蒙德恒等式\index{范德蒙德恒等式}。

取 $m = n = k$：$\binom{2n}{n} = \sum_{j=0}^{n} \binom nj
\binom{n}{n-j} = \sum_{j=0}^{n} \binom nj^2$，利用 $\binom{n}{n-j} =
\binom nj$。
\end{solution}

\begin{solution}{exo:b1:counting:8}
设 $A_d$ 是 $d$ 的倍数组成的\omterm{def:b1:logic:sets}{集合}（取自 $\intint{1}{1000}$），于是
$\abs{A_d} = \lfloor 1000/d \rfloor$。由容斥原理
（见\cref{thm:b1:counting:inclexcl}），取 $A_2, A_3, A_5$，并注意
$A_2 \cap A_3 = A_6$ 等：
\[
500 + 333 + 200 - 166 - 100 - 66 + 33 = 734 .
\]
故有 $734$ 个整数可被 $2$、$3$ 或 $5$ 整除。
\end{solution}

\begin{solution}{exo:b1:counting:9}
\omterm{def:b1:logic:inj}{满射}到 $2$ 个元素上：全部 $2^4 = 16$ 个\omterm{def:b1:logic:map}{映射}去掉 $2$ 个常值\omterm{def:b1:logic:map}{映射}，
共 $14$ 个\omterm{def:b1:logic:inj}{满射}。

\omterm{def:b1:logic:inj}{满射}到 $3$ 个元素上：对漏掉的值用容斥原理，
从 $4$ 元集到 $3$ 元集、至少漏掉一个值的\omterm{def:b1:logic:map}{映射}个数为
$\binom 31 2^4 - \binom 32 1^4 = 48 - 3 = 45$；全部\omterm{def:b1:logic:map}{映射}共 $3^4 =
81$，\omterm{def:b1:logic:inj}{满射}个数为 $81 - 45 = 36$。（检验：从 $4$ 元集到 $3$ 元集的\omterm{def:b1:logic:inj}{满射}
恰使一个值被取两次：选定被取两次的值（$3$ 种），
取到该值的两个元素（$\binom 42 = 6$ 种），以及其余部分的\omterm{def:b1:logic:inj}{双射}（$2$ 种）：
$3 \times 6 \times 2 = 36$。）
\end{solution}

\begin{solution}{exo:b1:counting:10}
方程 $x_1 + \dots + x_n = k$ 在 $\N^n$ 中的解，可编码为
由 $k$ 个星号与 $n - 1$ 个竖线组成的一行：依次写 $x_1$ 个星号、一个竖线、$x_2$ 个星号、
一个竖线、\dots，最后以 $x_n$ 个星号结束。这是
到长为 $k + n - 1$、恰含 $k$ 个星号与 $n - 1$ 个竖线的字上的\omterm{def:b1:logic:inj}{双射}，
而这些字由星号所在的位置唯一确定：
$\binom{n + k - 1}{k}$。可重选取对应于
方程的解（$x_i$ 为对象 $i$ 被选取的次数），
因此计数结果相同。
\end{solution}

\begin{solution}{exo:b1:counting:11}
记 $A_i = \{\sigma : \sigma(i) = i\}$，则
$\bigcap_{i \in I} A_i$ 中的\omterm{def:b1:counting:objects}{置换}固定每个 $i \in I$，并自由地\omterm{def:b1:counting:objects}{置换}其余的
$n - \abs I$ 个点，故 $\abs{\bigcap_{i \in I} A_i} =
(n - \abs I)!$。由容斥原理：
\[
\Bigl|\bigcup_i A_i\Bigr|
= \sum_{k=1}^{n} (-1)^{k+1} \binom nk (n-k)!
= \sum_{k=1}^{n} (-1)^{k+1} \frac{n!}{k!} ,
\]
由于共有 $\binom nk$ 个子集 $I$，大小均为 $k$，于是
\[
D_n = n! - \Bigl|\bigcup_i A_i\Bigr|
= n!\Bigl(1 - \sum_{k=1}^{n} \frac{(-1)^{k+1}}{k!}\Bigr)
= n! \sum_{k=0}^{n} \frac{(-1)^k}{k!} .
\]
对第二个等式，设 $\sigma$ 为 $\intint{1}{n}$ 上的\omterm{def:b1:counting:objects}{置换}，
按其不动点集 $F(\sigma)$ 分类。对一个固定的
$k$ 元子集 $F$，满足 $F(\sigma) = F$ 的\omterm{def:b1:counting:objects}{置换}恰为
补集上的错排，共 $D_{n-k}$ 个。
对 $\binom nk$ 种 $F$ 的选择，按每个 $k$ 求和：
$n! = \sum_{k=0}^{n} \binom nk D_{n-k}$.
\end{solution}

\begin{solution}{exo:b1:counting:12}
\emph{第一个恒等式。}计数对 $(A, a)$，其中 $A \subseteq E$
（$\abs E = n$），$a \in A$。按 $A$ 的大小，共 $\sum_k \binom nk k$
个对。先选标记元素：有 $n$ 种选法给 $a$，然后
再从其余 $n - 1$ 个元素中任取子集以补全 $A$：
$n\,2^{n-1}$ 个对。

\emph{第二个恒等式。}计数三元组 $(A, a, b)$，使 $a, b \in
A$（可能 $a = b$）。按大小计数：$\sum_k k^2 \binom nk$。直接计算：
要么 $a = b$（$n\,2^{n-1}$ 个三元组，即前面的计数），要么 $a \neq b$
（$n(n-1)$ 种有序选取，再从其余 $n - 2$
个元素中任取子集：$n(n-1)\,2^{n-2}$ 个）。总计
\[
n\,2^{n-1} + n(n-1)\,2^{n-2} = n\,2^{n-2}\,(2 + n - 1)
= n(n+1)\,2^{n-2} .
\]
\end{solution}

\begin{solution}{pb:b1:counting:1}
\textbf{1.} $D_1 = 0$（唯一的\omterm{def:b1:counting:objects}{置换}固定 $1$），$D_2 = 1$
（对换），$D_3 = 2$（单行记号下为 $231$ 与 $312$）。
对 $n = 4$，按 $\sigma(1)$ 分组：$\sigma(1) = 2$ 时的错排为
$2143$、$2341$、$2413$；$\sigma(1) = 3$ 时为 $3142$、$3412$、
$3421$；$\sigma(1) = 4$ 时为 $4123$、$4312$、$4321$。每组
各三个，故 $D_4 = 9$。

\textbf{2.} 恰有 $k$ 个不动点的\omterm{def:b1:counting:objects}{置换}，
由其不动点集 $F$ 的选取（$\binom nk$
种）与其在补集上的限制共同确定，该限制
必须是 $n - k$ 个点上的、\emph{没有}不动点的\omterm{def:b1:counting:objects}{置换}
（$D_{n-k}$ 种）。两种选择相互独立，且
其对应关系为\omterm{def:b1:logic:inj}{双射}：$P_k(n) = \binom nk D_{n-k}$。

\textbf{3.} $P_0(4) = D_4 = 9$；$P_1(4) = \binom41 D_3 = 4 \times 2
= 8$；$P_2(4) = \binom42 D_2 = 6$；$P_3(4) = \binom43 D_1 = 0$
（三个不动点必迫使第四个也为不动点）；$P_4(4) = 1$。和为 $9 + 8 + 6
+ 0 + 1 = 24 = 4!$。无匹配（$9$ 种）
胜过恰有一个匹配（$8$ 种）——仅以微弱优势。

\textbf{4.} 计数对 $(\sigma, i)$，满足 $\sigma(i) = i$。
对固定的 $i$，固定 $i$ 的\omterm{def:b1:counting:objects}{置换}即为
其余 $n - 1$ 个点上的\omterm{def:b1:counting:objects}{置换}，共 $(n-1)!$ 个。因此对的个数为
$n \cdot (n-1)! = n!$，此数也等于
$\sum_\sigma \abs{\mathrm{Fix}(\sigma)}$。除以\omterm{def:b1:counting:objects}{置换}个数
$n!$ 个\omterm{def:b1:counting:objects}{置换}，平均不动点个数为
恰为 $1$，对每个 $n \geq 1$ 成立。

\textbf{5.} 设 $\sigma$ 为 $\intint1{n+1}$ 的错排，
令 $j = \sigma(n+1) \in \intint1n$，共 $n$ 种可能取值。\emph{情形
$\sigma(j) = n+1$：}点 $j$ 与 $n+1$ 互换，而 $\sigma$
在其余 $n - 1$ 个点上的限制为任一错排，
共 $D_{n-1}$ 种可能。\emph{情形 $\sigma(j)
\neq n+1$：}设 $i_0 = \sigma^{-1}(n+1)$，此时 $i_0 \neq j$ 且
$i_0 \leq n$。定义 $\tau$ 为 $\intint1n$ 上的\omterm{def:b1:logic:map}{映射}，使 $\tau(i) = \sigma(i)$
对 $i \neq i_0$ 成立，且 $\tau(i_0) = j$。于是 $\tau$ 为
$\intint1n$ 上的\omterm{def:b1:counting:objects}{置换}（$n+1$ 被缺失的
值 $j$ 替代），且它为错排：$\tau(i_0) = j \neq i_0$，且
$\tau(i) = \sigma(i) \neq i$ 在其余处成立。反之，
由 $\tau$（它是 $\intint1n$ 的错排）与值 $j$，可还原
$\sigma$：令 $\sigma(n+1) = j$，$\sigma(\tau^{-1}(j)) = n+1$，
且其余处 $\sigma = \tau$。这是\omterm{def:b1:logic:inj}{双射}，共给出 $D_n$
种可能。对 $j$ 求和，得 $D_{n+1} = n(D_n + D_{n-1})$。
数值上：$D_5 = 4(9 + 2) = 44$，$D_6 = 5(44 + 9) = 265$。

\textbf{6.} 由第 5 题，$D_{n+1} = nD_n + nD_{n-1}$，于是
\[
u_{n+1} = D_{n+1} - (n+1)D_n = nD_n + nD_{n-1} - (n+1)D_n
= -(D_n - nD_{n-1}) = -u_n .
\]
由 $u_1 = D_1 - 1 \cdot D_0 = -1$，归纳可得 $u_n =
(-1)^n$，即 $D_n = nD_{n-1} + (-1)^n$，对 $n \geq 1$ 成立。

\textbf{7.} 对 $n$ 归纳。基础：$D_0 = 1 = 0!\,s_0$。归纳步：
假设 $D_{n-1} = (n-1)!\,s_{n-1}$，
\[
D_n = nD_{n-1} + (-1)^n = n!\,s_{n-1} + (-1)^n
= n!\Bigl(s_{n-1} + \frac{(-1)^n}{n!}\Bigr) = n!\,s_n ,
\]
此即所求公式。全程未用容斥原理，只用到
第 5 题中的组合递推关系。

\textbf{8.} 三项式修订，用阶乘化简：
\[
\binom nk \binom kj
= \frac{n!}{k!\,(n-k)!} \cdot \frac{k!}{j!\,(k-j)!}
= \frac{n!}{j!\,(n-j)!} \cdot \frac{(n-j)!}{(k-j)!\,(n-k)!}
= \binom nj \binom{n-j}{k-j} .
\]
代入 $a_k = \sum_j \binom kj b_j$，并交换两个
有限和的求和次序：
\[
\sum_{k=0}^{n} (-1)^{n-k} \binom nk a_k
= \sum_{j=0}^{n} b_j \binom nj
\sum_{k=j}^{n} (-1)^{n-k} \binom{n-j}{k-j}
= \sum_{j=0}^{n} b_j \binom nj
\sum_{i=0}^{n-j} (-1)^{(n-j)-i} \binom{n-j}{i} .
\]
内层和即 $(1 + (-1))^{n-j} = 0^{n-j}$ 的展开式
（二项式定理，\cref{thm:b1:counting:binomial}）：它
在 $j < n$ 时为零，等于 $1$ 的则是 $j = n$ 的情形。只有 $j = n$ 项保留，
右端即为 $b_n$，如所证。

\textbf{9.} 由对称性 $\binom nk = \binom n{n-k}$，
\cref{exo:b1:counting:11}中的恒等式可改写为 $n! =
\sum_{k=0}^n \binom nk D_k$。应用第 8 题，取 $a_n = n!$，
$b_k = D_k$：
\[
D_n = \sum_{k=0}^{n} (-1)^{n-k} \binom nk k!
= \sum_{k=0}^{n} (-1)^{n-k} \frac{n!}{(n-k)!}
= n! \sum_{j=0}^{n} \frac{(-1)^j}{j!} ,
\]
换指标 $j = n - k$：这是第三次用到该公式。

\textbf{10.} $D_n = n!\,s_n$（第 7 题），于是
\[
\Bigl| D_n - \frac{n!}{\eu} \Bigr|
= n!\,\abs{s_n - \eu^{-1}} < \frac{n!}{(n+1)!} = \frac1{n+1} .
\]

\textbf{11.} 当 $n \geq 1$ 时，$\frac1{n+1} \leq \frac12$，且
第 10 题的不等式是严格的：$D_n$ 与某点的距离
$< \frac12$，此点即 $n!/\eu$，故它是唯一的最近整数。
当 $n = 0$ 时，估计只给出距离 $< 1$，事实上结论
在此处不成立：$0!/\eu \approx 0.368$ 的最近整数为 $0$，而
$D_0 = 1$。

\textbf{12.} $\eu^{-1} - s_n = \sum_{k \geq n+1} (-1)^k/k!$ 是一个
各项严格递减的交错级数，故其符号为
首项 $(-1)^{n+1}/(n+1)!$ 的符号。因此 $s_n -
\eu^{-1}$ 的符号与 $(-1)^n$ 相同：$n$ 为偶数时，$s_n > \eu^{-1}$
从而 $D_n = n!\,s_n > n!/\eu$；$n$ 为奇数时，$D_n < n!/\eu$。

\textbf{13.} $D_7 = 6(265 + 44) = 6 \times 309 = 1854$；$D_8 =
7(1854 + 265) = 7 \times 2119 = 14\,833$；$D_9 = 8(14\,833 + 1854)
= 8 \times 16\,687 = 133\,496$；$D_{10} = 9(133\,496 + 14\,833) =
9 \times 148\,329 = 1\,334\,961$。验算：$10!/\eu = 3\,628\,800 /
2.718281828 \approx 1\,334\,960.92$，其最近整数为
$1\,334\,961$——且 $D_{10} > 10!/\eu$，与第 12 题
对偶数 $n$ 的预测相符。

\textbf{14.} $\abs{p_n - \eu^{-1}} = \abs{s_n - \eu^{-1}} <
\frac1{(n+1)!}$。当 $n = 6$ 时：$p_6 = 265/720 = 0.36806$（取五位小数），
与 $\eu^{-1} = 0.36788$ 相比，二者之差小于 $1/7! =
1/5040 < 2 \times 10^{-4}$。估计式 $1/(n+1)!$ 衰减得极快，
以至于只需十二封信，概率便已在许多位小数上
确定，因此答案“约为 $36.8\%$”，对一切
实际用途而言都与 $n$ 无关——这正是本题著名的
出人意料之处。

\textbf{15.} 由第 2 题及 $D_m = m!\,s_m$：
\[
\frac{P_k(n)}{n!} = \frac{\binom nk D_{n-k}}{n!}
= \frac{D_{n-k}}{k!\,(n-k)!} = \frac{s_{n-k}}{k!}
\;\longrightarrow\; \frac{\eu^{-1}}{k!}
\]
当 $n \to \infty$ 且 $k$ 固定时，由于 $s_{n-k} \to \eu^{-1}$。
极限值 $\eu^{-1}/k!$（$k \in \N$）是
参数为 $1$ 的泊松分布的权重。

\textbf{16.} 计数三元组 $(\sigma, i, j)$，其中 $i \neq j$，
$\sigma(i) = i$，$\sigma(j) = j$。先选定有序对，
共 $n(n-1)$ 种方式；同时固定 $i$ 与 $j$ 的\omterm{def:b1:counting:objects}{置换}
即其余 $n - 2$ 个点上的\omterm{def:b1:counting:objects}{置换}，共 $(n-2)!$ 个。
总数为 $n(n-1)(n-2)! = n!$。若改为先对 $\sigma$ 求和，
则对每个 $\sigma$，计数的是其相异不动点的有序对，
即 $\abs{\mathrm{Fix}(\sigma)}(\abs{\mathrm{Fix}(\sigma)}-1)$。
由此得所述恒等式。两边除以 $n!$，则
$\abs{\mathrm{Fix}}(\abs{\mathrm{Fix}} - 1)$ 的平均值为 $1$，故
$\abs{\mathrm{Fix}}^2$ 的平均值为 $1 + 1 = 2$，而方差为 $2 -
1^2 = 1$。

\textbf{17.} 比例 $1 - p_n$：当 $n = 4$ 时，$1 - \frac
9{24} = \frac{15}{24} = 0.6250$；当 $n = 5$ 时，$1 - \frac{44}{120} =
\frac{76}{120} = 0.6333$；当 $n = 6$ 时，$1 - \frac{265}{720} =
\frac{455}{720} = 0.6319$。各值与 $1 - \eu^{-1}
\approx 0.6321$ 相差不到百分之一，并在其附近上下振荡。

\textbf{18.} 直接计算：
\[
s_{n+2} - s_n = \frac{(-1)^{n+1}}{(n+1)!} +
\frac{(-1)^{n+2}}{(n+2)!}
= (-1)^{n+1}\Bigl(\frac1{(n+1)!} - \frac1{(n+2)!}\Bigr),
\]
括号内的量 $> 0$。当 $n$ 为偶数时，此差
为负：$s_{n+2} < s_n$，故 $p_0 > p_2 > p_4 > \dots$；当 $n$
为奇数时则为正，即 $p_1 < p_3 < p_5 < \dots$。结合
第 12 题（偶数项在 $\eu^{-1}$ 之上，奇数项在其下）与
第 14 题（与 $\eu^{-1}$ 的距离趋于 $0$）：两条阶梯状数列
将 $\eu^{-1}$ 夹在中间。

\textbf{19.} 一次完整抽签即一个均匀随机\omterm{def:b1:counting:objects}{置换}，
当它是错排时有效，其概率为 $p_n \approx \eu^{-1}$。
由所引事实，直至成功所需的平均抽签次数为
$1/p_n$，由第 14 题得 $1/p_n \approx \eu$，误差
在 $n$ 较小时即可忽略。因此，“秘密圣诞老人”
带重抽的游戏平均约需 $\eu \approx 2.72$ 次完整抽签
——无论办公室有 $6$ 人还是 $600$ 人。

\textbf{20.} 固定 $j \geq 2$，对
取值 $\sigma(1) = j$ 运用第 5 题的分类。若 $\sigma(j) = 1$，则其余
$n - 2$ 个点上可取任一错排，共 $D_{n-2}$ 种。若
$\sigma(j) \neq 1$：将\omterm{def:b1:logic:map}{原像} $i_0 = \sigma^{-1}(1)$ 改指向
$j$，与第 5 题相同。这是一个\omterm{def:b1:logic:inj}{双射}，
与 $n - 1$ 个点 $\{2, \dots, n\}$ 上的错排一一对应，共 $D_{n-1}$
种。总数为 $D_{n-1} + D_{n-2}$，对每个 $j$ 均相同。对
$n - 1$ 个 $j$ 求和，得 $D_n = (n-1)(D_{n-1} + D_{n-2})$，
其中显含因子 $n - 1$，即 $(n-1) \mid D_n$。

\textbf{21.} 断言：$D_n$ 为奇数当且仅当 $n$ 为偶数。用归纳法，
利用 $D_n = nD_{n-1} + (-1)^n$，即 $D_n \equiv nD_{n-1} + 1 \pmod 2$。
奠基：$D_1 = 0$ 为偶数，$n = 1$ 为奇数，断言成立。若 $n$ 为偶数，
则 $nD_{n-1}$ 为偶数，于是 $D_n \equiv 1$ 为奇数，与断言相符。若 $n$ 为
奇数，则 $n - 1$ 为偶数，由归纳假设 $D_{n-1}$ 为奇数，且
$D_n \equiv D_{n-1} + 1 \equiv 0$，为偶数。归纳证明至此完成。

\textbf{22.} 将 $D_n = nD_{n-1} + (-1)^n$ 模 $n$ 约化，首项
消去，得 $D_n \equiv (-1)^n \pmod n$。取 $n = 10$，
则 $(-1)^{10} = 1$，且 $D_{10} = 1\,334\,961$ 的末位数字
为 $1$，与之相符。

\textbf{23.} 当 $n \geq 3$ 时，$D_{n-1} \geq 1$，于是对第 6 题的
递推式除以 $D_{n-1}$，得 $D_n/D_{n-1} = n +
(-1)^n/D_{n-1}$，且 $\abs{(-1)^n/D_{n-1}} \leq 1$，并迅速
趋于 $0$。一致性：若 $D_n \approx n!/\eu$，则
$D_n/D_{n-1} \approx n!/(n-1)! = n$——因子 $\eu$ 在比值中约去，
递推式可确认这一点，精度为
$1/D_{n-1}$。

\textbf{24.} (i) 乘法法则与加法法则是一切计数的基础：
第 2 题与第 5 题把\omterm{def:b1:counting:objects}{置换}的\omterm{def:b1:logic:sets}{集合}\omterm{thm:b1:logic:partition}{划分}为相互独立的
若干阶段。(ii) 双重计数求出了不动点个数的均值（第 4 题）与
方差（第 16 题），且完全不需要 $D_n$ 的任何公式。
(iii) 二项式定理计算了
使二项式反演成立的交错内和 $(1-1)^{n-j}$
（第 8 题）。(iv) 交错级数估计把
精确却晦涩的和 $n!\,s_n$ 化为直观的结论
“最接近 $n!/\eu$ 的整数”（第 10--14 题）。

\textbf{25.} 容斥原理
（\cref{ex:b1:counting:derangement}与\cref{exo:b1:counting:11}）
是概念性的证明，它把交错和解释为
多计数的修正，并且可以原样推广到计数
避开任一族“坏”\omterm{def:b1:logic:sets}{集合}的元素。递推路线
（第 5--7 题）计算最快——线性时间、精确整数运算、
无需阶乘——并且是
第五部分算术事实的来源。二项式反演（第 8--9 题）把
该公式置于一般的变换之中，此类变换将在
任何两个相互对立的三角恒等式组出现处再次出现。而
$\eu$ 的出现最好由公式本身来解释：
错排所占比例是 $s_n$，它是 $\eu^{-1}$ 级数的部分和，
因此蒙特莫尔的信封问题早在欧拉采用这一记号的三十年前，
便已在计算数 $\eu$ 了。
\end{solution}
